\begin{figure}[htbp]
\begin{center}
\includegraphics[width=1.0\textwidth]{FlxGrp_2x2_20261026.pdf}
\end{center}
\caption{Biomimetic-knee torque versus knee angle for the flexor BPAs, in a 2$\times$2 layout with the legend in the fourth tile. Knee angle is negative in flexion, and flexion torque is negative. \textbf{(A)} 10 mm panel: $\diameter$\qty{10}{\mm}, $l_{0}=\qty{41.5}{\cm}$ actuator with a \qty{21}{\mm} artificial tendon at \qty{604}{\kPa}. \textbf{(B)} 20 mm panel: $\diameter$\qty{20}{\mm}, $l_{0}=\qty{41.5}{\cm}$ actuator with a \qty{14}{\mm} artificial tendon at \qty{620}{\kPa}. \textbf{(C)} 20 mm panel, the redesigned two-BPA flexor route (two $\diameter$\qty{20}{\mm} BPAs, $l_{0}=\qty{56.2}{\cm}$) at \qty{620}{\kPa}, a different configuration from panel \textbf{(B)} despite the shared title: corrected prediction and measured points against the human target. Indigo filled circles are measured torque; the magenta line is the optimized prediction of the corrected model using the flexor terms identified on the pinned knee; the dashed light orange line is the unoptimized prediction of the improved BPA characterization alone. In \textbf{(B)}, the black dotted line is the maximum human isometric torque, and the pink circles are measured points at the reduced pressures \qtylist{560;421;325;281}{\kPa} (from left to right, respectively), recorded because full pressure would load the printed femur beyond its ultimate strength; in \textbf{(C)}, the black dotted line is the maximum human isometric torque of the biceps femoris short head, and across its fifteen tested data measurements from \ang{-119} to \ang{5} the measured torque magnitude met or exceeded the human target at every point, with margins from \qty{0.42}{\newton\meter} at \ang{-105} to \qty{17.0}{\newton\meter} at \ang{-14}. Panels \textbf{(B)} and \textbf{(C)} share their axis ranges.}\label{fig:FlxGrp}
\end{figure}
